% Bagian: intuitionistic-logic
% Bab: introduction
% Subbagian: axiomatic-derivations

\documentclass[../../../include/open-logic-section]{subfiles}

\begin{document}

\olfileid[id]{int}{int}{axd}

\olsection{\usetoken{P}{derivation} Aksiomatik}

!!{derivation}s aksiomatik bagi logika proposisional intuisionistik merupakan
sistem !!{derivation} yang paling sederhana secara konseptual dan yang pertama
secara historis. Sistem ini bekerja tepat seperti dalam logika proposisional
klasik.

\begin{defn}[!!^{derivability}]
Jika $\Gamma$ merupakan himpunan !!{formula}s dari $\Lang L$, maka suatu
\emph{!!{derivation}} dari $\Gamma$ ialah urutan berhingga $!A_1$, \dots,
~$!A_n$ dari !!{formula}s sedemikian sehingga, bagi setiap $i \le n$, salah
satu dari hal berikut berlaku:
\begin{enumerate}
\item $!A_i \in \Gamma$; atau
\item $!A_i$ merupakan suatu aksioma; atau
\item $!A_i$ diperoleh dari suatu $!A_j$ dan $!A_k$, dengan $j < i$ dan $k <
  i$, melalui modus ponens; yakni, $!A_k \ident !A_j \lif !A_i$.
\end{enumerate}
\end{defn}

\begin{defn}[Aksioma]
Himpunan~$\PAx$ dari \emph{aksioma} bagi logika proposisional intuisionistik
terdiri atas semua !!{formula}s dalam bentuk-bentuk berikut:
\begin{align}
  & (!A \land !B) \lif !A \ollabel{ax:land1}\\
  & (!A \land !B) \lif !B \ollabel{ax:land2}\\
  & !A \lif (!B \lif (!A \land !B)) \ollabel{ax:land3}\\
  & !A \lif (!A \lor !B) \ollabel{ax:lor1}\\
  & !A \lif (!B \lor !A) \ollabel{ax:lor2}\\
  & (!A \lif !C) \lif ((!B \lif !C) \lif ((!A \lor !B) \lif !C)) \ollabel{ax:lor3}\\
  & !A \lif (!B \lif !A) \ollabel{ax:lif1}\\
  & (!A \lif (!B \lif !C)) \lif ((!A \lif !B) \lif (!A \lif !C)) \ollabel{ax:lif2}\\
  & \lfalse \lif !A \ollabel{ax:lfalse1}
\end{align}
\end{defn}

\begin{defn}[!!^{derivability}]
!!^a{formula}~$!A$ \emph{!!{derivable}} dari $\Gamma$, ditulis $\Gamma
\Proves !A$, jika terdapat !!a{derivation} dari $\Gamma$ yang berakhir pada
$!A$.
\end{defn}

\begin{defn}[Teorema]
!!^a{formula}~$!A$ merupakan \emph{teorema} jika terdapat !!a{derivation}
bagi~$!A$ dari himpunan kosong. Kita menulis $\Proves !A$ jika $!A$ merupakan
teorema, dan $\Proves/ !A$ jika bukan.
\end{defn}

\begin{prop}
  Jika $\Gamma \Proves !A$ dalam sistem !!{derivation} aksiomatik bagi logika
  intuisionistik, maka $\Gamma \Proves !A$ juga dalam logika klasik. Khususnya,
  jika $!A$ merupakan teorema intuisionistik, formula itu juga merupakan
  teorema klasik.
\end{prop}

\begin{proof}
  Setiap aksioma intuisionistik juga merupakan aksioma klasik, sehingga setiap
  !!{derivation} dalam logika intuisionistik juga merupakan !!a{derivation}
  dalam logika klasik.
\end{proof}

\end{document}
